% ============================================================
\chapter{Population Dynamics Models}
\label{ch:populations}
% ============================================================

\section{Introduction}

In 1798, the Reverend Thomas Malthus published his famous \emph{Essay on the Principle of Population}, arguing that population grows geometrically while resources grow only arithmetically. This pessimistic vision---the ``Malthusian catastrophe''---would inspire Darwin and become the first mathematical model of population dynamics: the exponential. But Malthus was wrong on one crucial point: growth cannot be exponential forever. In 1838, Pierre-Fran\c{c}ois Verhulst introduced a saturation term leading to the celebrated logistic equation. Then, in the 1920s, Alfred Lotka and Vito Volterra, working independently---one on chemical reactions, the other on Adriatic fisheries---discovered the predator-prey model that bears their names.

This chapter retraces this intellectual adventure and develops the fundamental models --- Malthus,
Verhulst, Lotka--Volterra --- emphasising derivation from first
principles, qualitative analysis (phase portraits, stability), and
numerical simulations.

% ------------------------------------------------------------
\section{Exponential Growth: The Malthus Model}
\label{sec:malthus}
% ------------------------------------------------------------

\begin{definition}[Malthus model]
Let $N(t)$ be the population size at time $t$.  If the birth rate $b$
and death rate $d$ are constant and proportional to population size:
\[
  \frac{dN}{dt} = rN, \qquad r = b - d, \qquad N(0) = N_0.
\]
\end{definition}

\begin{theorem}[Solution of the Malthus model]
The solution is $N(t) = N_0\,e^{rt}$.
\begin{itemize}
  \item $r > 0$: exponential growth (explosion).
  \item $r = 0$: constant population.
  \item $r < 0$: exponential decay (extinction).
\end{itemize}
\end{theorem}

\begin{remark}
The Malthus model is realistic only over short periods.  No real
population can grow exponentially indefinitely.
\end{remark}

\begin{codeblock}[title=\textbf{Malthus model in Python}]
\begin{minted}{python}
import numpy as np
import matplotlib.pyplot as plt

t = np.linspace(0, 10, 200)
N0 = 100

for r in [-0.3, 0, 0.2, 0.5]:
    N = N0 * np.exp(r * t)
    plt.plot(t, N, label=f'$r = {r}$')

plt.xlabel('Time $t$'); plt.ylabel('Population $N(t)$')
plt.legend(); plt.title('Malthus Model')
plt.ylim(0, 2000)
plt.tight_layout(); plt.savefig('malthus.pdf')
\end{minted}
\end{codeblock}

% ------------------------------------------------------------
\section{Logistic Growth: The Verhulst Model}
\label{sec:verhulst}
% ------------------------------------------------------------

\begin{definition}[Verhulst logistic model]
To model resource limitation, Verhulst (1838) proposed:
\[
  \frac{dN}{dt} = rN\left(1 - \frac{N}{K}\right),
  \qquad N(0) = N_0,
\]
where $K > 0$ is the \textbf{carrying capacity}.
\end{definition}

\begin{intuition}[title=\textbf{Interpreting the logistic term}]
The factor $(1 - N/K)$ models intraspecific competition: when
$N \ll K$, growth is nearly exponential; as $N$ approaches $K$, the
effective growth rate tends to zero.
\end{intuition}

\begin{theorem}[Solution of the logistic equation]
The logistic ODE is separable.  Its solution is:
\[
  N(t) = \frac{K}{1 + \left(\frac{K}{N_0}-1\right)e^{-rt}}.
\]
We have $\lim_{t\to+\infty} N(t) = K$ for all $N_0 > 0$.
\end{theorem}

\begin{proof}
Separate variables:
\[
  \frac{dN}{N(1-N/K)} = r\,dt.
\]
Partial fractions:
\[
  \frac{1}{N(1-N/K)} = \frac{1}{N} + \frac{1/K}{1-N/K}.
\]
Integration:
\[
  \ln N - \ln(1-N/K) = rt + C.
\]
Setting $C = \ln\frac{N_0}{1-N_0/K}$ yields the result.
\end{proof}

\subsection{Qualitative Analysis}

Equilibria solve $rN(1-N/K)=0$: $N^*=0$ and $N^*=K$.

\begin{proposition}[Stability of logistic equilibria]
\begin{itemize}
  \item $N^*=0$ is \textbf{unstable} (for $r>0$).
  \item $N^*=K$ is \textbf{asymptotically stable}.
\end{itemize}
This follows from linearisation: $f'(N)=r(1-2N/K)$, giving
$f'(0)=r>0$ (unstable) and $f'(K)=-r<0$ (stable).
\end{proposition}

\begin{codeblock}[title=\textbf{Logistic growth and phase portrait}]
\begin{minted}{python}
import numpy as np
import matplotlib.pyplot as plt
from scipy.integrate import odeint

r, K = 0.5, 1000

def logistic(N, t):
    return r * N * (1 - N / K)

t = np.linspace(0, 30, 500)

fig, (ax1, ax2) = plt.subplots(1, 2, figsize=(12, 4))

for N0 in [10, 50, 200, 500, 1200, 1500]:
    N = odeint(logistic, N0, t).flatten()
    ax1.plot(t, N, label=f'$N_0={N0}$')

ax1.axhline(K, color='gray', ls='--', label=f'$K={K}$')
ax1.set_xlabel('$t$'); ax1.set_ylabel('$N(t)$')
ax1.legend(fontsize=8); ax1.set_title('Trajectories')

N_vals = np.linspace(0, 1600, 300)
dNdt = r * N_vals * (1 - N_vals / K)
ax2.plot(N_vals, dNdt, 'b-')
ax2.axhline(0, color='gray', ls='--')
ax2.plot([0, K], [0, 0], 'ro', markersize=8)
ax2.set_xlabel('$N$'); ax2.set_ylabel("$dN/dt$")
ax2.set_title('Velocity field')
plt.tight_layout(); plt.savefig('logistic.pdf')
\end{minted}
\end{codeblock}

% ------------------------------------------------------------
\section{The Lotka--Volterra Predator--Prey Model}
\label{sec:lotka_volterra}
% ------------------------------------------------------------

\begin{definition}[Lotka--Volterra system]
Let $x(t)$ be the prey population and $y(t)$ the predator population.
The classical Lotka--Volterra system is:
\begin{align}
  \dot{x} &= ax - bxy, \label{eq:lv_prey} \\
  \dot{y} &= -cy + dxy, \label{eq:lv_pred}
\end{align}
where $a,b,c,d>0$.
\begin{itemize}
  \item $a$: intrinsic prey growth rate.
  \item $b$: predation rate.
  \item $c$: predator death rate.
  \item $d$: conversion efficiency prey $\to$ predator.
\end{itemize}
\end{definition}

\subsection{Equilibria}

Setting the right-hand side to zero:
\[
  x(a-by)=0, \qquad y(-c+dx)=0.
\]

\begin{proposition}[Lotka--Volterra equilibria]
There are two equilibria:
\begin{enumerate}
  \item $E_0 = (0,0)$: total extinction.
  \item $E^* = (c/d,\;a/b)$: coexistence.
\end{enumerate}
\end{proposition}

\subsection{Linear Stability}

The Jacobian is:
\[
  J(x,y) = \begin{pmatrix}
    a-by & -bx \\ dy & -c+dx
  \end{pmatrix}.
\]

\begin{theorem}[Stability of Lotka--Volterra equilibria]
\begin{enumerate}
  \item At $E_0=(0,0)$: $J = \begin{pmatrix}a&0\\0&-c\end{pmatrix}$,
    eigenvalues $\lambda_1=a>0$, $\lambda_2=-c<0$.
    $E_0$ is a \textbf{saddle point} (unstable).
  \item At $E^*=(c/d,a/b)$: $J = \begin{pmatrix}0&-bc/d\\da/b&0\end{pmatrix}$,
    eigenvalues $\lambda=\pm i\sqrt{ac}$.
    $E^*$ is a \textbf{centre} (periodic orbits).
\end{enumerate}
\end{theorem}

\subsection{First Integral and Closed Trajectories}

\begin{theorem}[First integral]
The Lotka--Volterra system admits the first integral:
\[
  H(x,y) = dx - c\ln x + by - a\ln y = \text{constant}.
\]
Phase-plane trajectories are level curves of $H$, forming closed
curves around $E^*$.
\end{theorem}

\begin{proof}
Compute $dH/dt$:
\begin{align*}
  \frac{dH}{dt} &= d\dot{x} - \frac{c}{x}\dot{x}
                  + b\dot{y} - \frac{a}{y}\dot{y} \\
  &= d(ax-bxy) - \frac{c}{x}(ax-bxy)
   + b(-cy+dxy) - \frac{a}{y}(-cy+dxy) \\
  &= dax - dbxy - ca + cby - bcy + bdxy + ac - adx = 0.
\end{align*}
\end{proof}

\begin{codeblock}[title=\textbf{Lotka--Volterra simulation}]
\begin{minted}{python}
import numpy as np
import matplotlib.pyplot as plt
from scipy.integrate import solve_ivp

a, b, c, d = 1.0, 0.1, 1.5, 0.075

def lotka_volterra(t, z):
    x, y = z
    return [a*x - b*x*y, -c*y + d*x*y]

t_span = (0, 40)
z0 = [40, 9]
sol = solve_ivp(lotka_volterra, t_span, z0,
                t_eval=np.linspace(0, 40, 1000),
                method='RK45', rtol=1e-10)

fig, (ax1, ax2) = plt.subplots(1, 2, figsize=(13, 5))

ax1.plot(sol.t, sol.y[0], 'b-', label='Prey $x(t)$')
ax1.plot(sol.t, sol.y[1], 'r-', label='Predators $y(t)$')
ax1.set_xlabel('$t$'); ax1.set_ylabel('Population')
ax1.legend(); ax1.set_title('Time series')

for x0 in [10, 20, 30, 40, 60]:
    sol_i = solve_ivp(lotka_volterra, (0, 50), [x0, 9],
                      t_eval=np.linspace(0, 50, 2000),
                      method='RK45', rtol=1e-10)
    ax2.plot(sol_i.y[0], sol_i.y[1], 'b-', alpha=0.6)

ax2.plot(c/d, a/b, 'ro', markersize=8, label="$E^*$")
ax2.set_xlabel('Prey $x$'); ax2.set_ylabel('Predators $y$')
ax2.legend(); ax2.set_title('Phase portrait')
plt.tight_layout(); plt.savefig('lotka_volterra.pdf')
\end{minted}
\end{codeblock}

% ------------------------------------------------------------
\section{Extensions of the Lotka--Volterra Model}
\label{sec:extensions}
% ------------------------------------------------------------

\subsection{Holling Functional Response}

\begin{definition}[Functional responses]
The functional response $g(x)$ models the per-predator predation rate
as a function of prey density:
\begin{itemize}
  \item \textbf{Type I (linear):} $g(x)=bx$ (classical Lotka--Volterra).
  \item \textbf{Type II (saturation):} $g(x)=\frac{bx}{1+bhx}$
    where $h$ is the handling time.
  \item \textbf{Type III (sigmoidal):} $g(x)=\frac{bx^2}{1+bhx^2}$.
\end{itemize}
\end{definition}

The Type II system (Rosenzweig--MacArthur model) is:
\begin{align}
  \dot{x} &= rx\left(1-\frac{x}{K}\right) - \frac{bxy}{1+bhx}, \\
  \dot{y} &= y\left(\frac{dbx}{1+bhx} - c\right).
\end{align}

\begin{warning}[title=\textbf{Paradox of enrichment}]
Rosenzweig (1971) showed that increasing the carrying capacity $K$
can destabilise the coexistence equilibrium through a Hopf
bifurcation, causing large-amplitude oscillations that may lead to
extinction.
\end{warning}

\subsection{Interspecific Competition}

\begin{definition}[Lotka--Volterra competition model]
For two competing species:
\begin{align}
  \dot{N}_1 &= r_1 N_1\left(1 - \frac{N_1 + \alpha_{12}N_2}{K_1}\right), \\
  \dot{N}_2 &= r_2 N_2\left(1 - \frac{N_2 + \alpha_{21}N_1}{K_2}\right),
\end{align}
where $\alpha_{ij}$ measures the competitive effect of species $j$ on
species $i$.
\end{definition}

\begin{proposition}[Competitive exclusion principle]
If $\alpha_{12}K_2/K_1 > 1$ and $\alpha_{21}K_1/K_2 > 1$, the two
species cannot stably coexist (Gause's competitive exclusion).  The
outcome depends on initial conditions.
\end{proposition}

% ------------------------------------------------------------
\section{Discrete Models: Beverton--Holt Equation}
% ------------------------------------------------------------

\begin{definition}[Beverton--Holt model]
For populations with non-overlapping generations:
\[
  N_{t+1} = \frac{RN_t}{1 + (R-1)N_t/K},
\]
where $R>1$ is the reproduction rate and $K$ the carrying capacity.
\end{definition}

\begin{proposition}
The nontrivial equilibrium $N^*=K$ is globally stable for $R>1$.
The Beverton--Holt model is the discrete analogue of the logistic
equation and does \emph{not} exhibit chaos, unlike the logistic map
$N_{t+1}=RN_t(1-N_t/K)$ studied in Chapter~\ref{ch:chaos}.
\end{proposition}

% ------------------------------------------------------------
\section{Limitations of Population Models}
% ------------------------------------------------------------

\begin{warning}[title=\textbf{Assumptions and limitations}]
\begin{itemize}
  \item \textbf{Spatial homogeneity:} the above models ignore spatial
    structure.  In reality, populations are distributed in space
    (see Chapter~\ref{ch:diffusion}).
  \item \textbf{Determinism:} for small populations, stochastic
    fluctuations become important (Chapter~\ref{ch:stochastic}).
  \item \textbf{Age structure:} Leslie models introduce age classes.
  \item \textbf{Constant parameters:} in reality, $r$ and $K$ may
    depend on time (seasonality, climate change).
\end{itemize}
\end{warning}

% ------------------------------------------------------------
\section{Exercises}
% ------------------------------------------------------------

\begin{exercise}
Show that the logistic population reaches half the carrying capacity at
time $t_{1/2} = \frac{1}{r}\ln\frac{K-N_0}{N_0}$.
\end{exercise}

\begin{exercise}
Add constant harvesting $h$ to the logistic model:
$\dot{N} = rN(1-N/K) - h$.
\begin{enumerate}
  \item Find the equilibria as a function of $h$.
  \item Determine the maximum sustainable yield $h_{\max}$.
  \item Draw the bifurcation diagram $N^*$ vs.\ $h$.
  \item Simulate in Python for $h < h_{\max}$, $h = h_{\max}$,
    $h > h_{\max}$.
\end{enumerate}
\end{exercise}

\begin{exercise}
For the Lotka--Volterra system, verify numerically that $H(x,y)$ is
conserved over time.  Compare results from \texttt{RK45} and
\texttt{RK23} in \texttt{solve\_ivp} for different tolerances.
\end{exercise}

\begin{exercise}
Study the Rosenzweig--MacArthur model with parameters $r=1$, $b=0.5$,
$h=0.5$, $d=0.5$, $c=0.3$ for $K\in\{5,10,20,50\}$.
\begin{enumerate}
  \item Find the equilibria and their stability.
  \item Draw phase portraits.
  \item Identify the critical value of $K$ for the Hopf bifurcation.
\end{enumerate}
\end{exercise}

\begin{exercise}
Two bacterial species compete in a confined environment with $r_1=0.8$,
$r_2=0.6$, $K_1=500$, $K_2=400$, $\alpha_{12}=0.7$, $\alpha_{21}=1.2$.
\begin{enumerate}
  \item Find the equilibria and analyse their stability.
  \item Simulate the system for different initial conditions.
  \item Conclude: is there coexistence or exclusion?
\end{enumerate}
\end{exercise}
